\subsection{FreeBSD}

FreeBSD is a \textbf{free and open-source Unix-like system} that is first conceptualized and developed by Nate Williams,
Rod Grimes, and Jordan Hubbard in 1993 \parencite{freebsd_intro}. FreeBSD is a direct descendant of \textit{386BSD}, which is
a free and redistributable Unix-like system designed by William "Bill" Jolitz and Lynne Jolitz released in 1992 \parencite{bentson1995}.
386BSD was designed specifically for Intel 80386-based IBM personal computers \parencites{mckusick1999}{bentson1995}.
It was also the first Unix-like system to be freely distributable on home computers \parencite{bentson1995}.
Shortly after, the development of 386BSD has slowed down, mainly because it has been neglected by Bill Jolitz for
over a year \parencite{mckusick1999}. This brought problems and bugs in the system that were not addressed by
the official patchkit system \parencite{freebsd_intro}. From this, the \textit{Unofficial 386BSD Patchkit} was formed \parencite{jude2020}
as an alternative to the official patchkit system. Eventually, there was a disagreement between the Jolitz and
the maintainers of the unofficial patchkit regarding the future directions of 386BSD \parencite{jolitz2014}.
This then led Williams, Grimes, and Hubbard, which were the remaining maintainers of Unofficial 386BSD Patchkit,
to develop their own iteration of 386BSD now known as FreeBSD \parencite{freebsd_intro}. FreeBSD was coined by David Greenman
in the 19th of June 1993 \parencite{greenman1993}. Other name suggestions that were brought up in the email sent by
\textcite{greenman1993} include \textit{"BSDFree86"} and \textit{"Free86BSD"}.

FreeBSD is mainly inspired by \textit{Unix}, which was originally developed by Ken Thompson and Dennis Ritchie
at AT\&T Bell Laboratories \parencite{urban_tiemman2001}. It was originally written in assembly and later rewritten in C,
which was designed by the two same developers \parencite{urban_tiemman2001}. Unix was based on \textit{Multics}, which
is an operating system designed for the GE-645 mainframe \parencite{freebsd_intro} which introduced concepts foundational in
the development of modern-day operating systems such as dynamic linking, hierarchical file system, and memory-mapped files
\parencite{freebsd_intro}. FreeBSD also carried the legacy of \textbf{Berkeley Software Distribution (BSD)}, founded by
University of California, Berkeley Computer Systems Research Group (CSRG) in the 1970s \parencite{freebsd_torchbearer}.
This legacy is carried by adopting the \textbf{BSD License}, which allows FreeBSD to be redistributed and/or modified
\parencite{opensource2011}. The distribution of BSD to various academic institutions also sparked the open-source
movement \parencite{jude2020}, which advocated for transparency in how software is developed.
The philosophy of open-source software has been rooted into FreeBSD ever since its direct ancestor, with the Jolitz
using the \textit{Network Release 2} or \textit{Net/2} as a base for building 386BSD. Net/2 originated from Net/1 or
Network Release 1, which implemented the OSI network protocol stack and several new TCP/IP algorithms for networking
\parencite{jude2020}. With the rising costs of AT\&T software licenses, several groups are interested in developing
systems without using AT\&T utilities. By the development of Net/2, it adopted barely any AT\&T code, which nearly made
it freely distributable \parencite{jude2020}. However, the maintainers of FreeBSD were forced to change their Net/2
base to a newer 4.4BSD Lite because of the Novell/Berkeley lawsuit \parencite{howard}. This lawsuit has impacted
the development of BSD, allowing Linux to flourish instead \parencite{hubbard_andrews2002}. Nevertheless, several
versions of BSD, which includes FreeBSD, were able to prosper and adapt in niche applications.

The FreeBSD Project aims to develop a system that excels in \textbf{performance, security, and stability}
\parencite{freebsd_faq}. According to \textcite{freebsd_intro}, the following highlights the particular
strengths of FreeBSD:

\begin{itemize}
	\item \textbf{Open-Source License}: As previously mentioned, FreeBSD is licensed under the BSD License, which
	      allows FreeBSD to be freely distributable with or without modifications. This allows FreeBSD to be adopted
	      and modified for server and embedded system applications. FreeBSD is also an excellent resource in researching
	      operating systems and other areas in computer science.

	\item \textbf{Industry-Standard TCP/IP Networking}: FreeBSD implements the industry-standard networking protocols
	      for TCP/IP, which makes it an ideal system for a wide range of networking services such as web servers, routing,
	      firewalls, FTP servers, email servers, storage servers, virtualization servers, etc.

	\item \textbf{OpenZFS Support}: FreeBSD natively supports OpenZFS, which is an open-source implementation of
	      the ZFS filesystem. The ZFS filesystem is notable for its efficient data compression and excellent data recovery
	      such as protection against data corruption, integrity checking, and snapshots \parencite{openzfs_project}.

	\item \textbf{Security}: FreeBSD limits the privileges assigned to services, which reduces impact that might
	      be caused by security breaches. FreeBSD also implements several security features such as the Mandatory Access
	      Control Framework, Capsicum capability, and sandboxing mechanisms.

	\item \textbf{Linux Compatibility}: A lot of packages and binaries that are designed and/or well-supported for Linux systems are
	      compatible with FreeBSD. This includes commonly used utilities such Curl, Nginx, and FFmpeg \parencite{freebsd_ch12};
	      development tools such as C/C++ compiler and debugging suite, Git, Apache, and SQL languages such as PostgreSQL and
	      MySQL \parencite{freebsd_ch12}; text editors such as Vim and Emacs \parencite{freebsd_ch12}; and desktop servers,
	      such as X11 and Wayland, and environments, such as GNOME and KDE.
\end{itemize}

With its BSD License, FreeBSD has been adopted and modified for a variety of server, desktop, and embedded systems.
FreeBSD has been utilized in ARM and AArch64-based embedded systems for routers, firewall, and appliances \parencite{freebsd_intro}.
Some BSD operating systems such as DragonFly BSD and GhostBSD, were derived from FreeBSD. Sony also modified FreeBSD
to design CeliOS and Orbis OS, which were the operating systems driving the Playstation 3 and 4 respectively
\parencites{sony2019}{larabel2013}. FreeBSD is also used to drive Network-Attached Storage (NAS) systems
with FreeNAS or now TrueNAS \parencite{posch2024}. A significant portion of FreeBSD's code was also borrowed by Apple
in developing the Darwin operating system, which has been the base system for modern Apple operating systems such
as MacOS and iOS. \parencite{apple2002}.
